\documentclass{../latex/utmreport}
\newcommand{\ReportTitle}{Lab 8: The Same Calls, Inside Your Own OS}
% Shared title-page and PDF metadata for every laboratory report.
\newcommand{\StudentName}{Gafenco Victor}
\newcommand{\StudentGroup}{FAF-241}
\newcommand{\CourseName}{Operating Systems}
\newcommand{\InstructorName}{Patricia Reițman}
\newcommand{\InstructorTitle}{Assistant lecturer}
\newcommand{\ReportCity}{Chișinău}
\newcommand{\ReportYear}{2026}
\newcommand{\CourseRepository}{https://github.com/yoda-dinmd/os-labs}

\selectlanguage{english}
\utmlogo{../latex/assets/UTM_Logo_English.jpeg}
\institution{Technical University of Moldova}
\faculty{Faculty of Computers, Informatics, and Microelectronics}
\department{Department of Software Engineering}
\studyprogramme{Software Engineering Study Programme}
\documenttype{Laboratory report}
\student{\StudentName}{\StudentGroup}
\supervisor{\InstructorName}{\InstructorTitle}
\renewcommand{\UTMCoordinatorLabel}{Instructor}
\cityyear{\ReportCity}{\ReportYear}
\reportsubtitle{\CourseName}
\hypersetup{pdfauthor={\StudentName},pdfsubject={\CourseName},
  pdftitle={\ReportTitle}}
\reporttitle{\ReportTitle}
\addto\captionsenglish{\renewcommand{\figurename}{Figure}\renewcommand{\tablename}{Table}\renewcommand{\bibname}{References}}
\renewcommand{\figurename}{Figure}
\renewcommand{\tablename}{Table}
\renewcommand{\bibname}{References}
\usepackage{float}
\usepackage{placeins}
\setlength{\emergencystretch}{2em}
\newcommand{\code}[1]{\texttt{\detokenize{#1}}}
% Keep screenshots beside their observations, with bounded height and width.
\newcommand{\screenshot}[4][\textwidth]{%
  \begin{figure}[H]
    \centering
    \includegraphics[width=#1,height=0.62\textheight,keepaspectratio]{screenshots/#2}
    \caption{#3}\label{#4}
  \end{figure}}
\newcommand{\reportreferences}[2]{%
  \clearpage\phantomsection\addcontentsline{toc}{chapter}{References}
  \begin{thebibliography}{9}
    \bibitem{handout} Technical University of Moldova, FCIM / FAF.
    \textit{#1}. Operating Systems laboratory handout, academic year 2026--2027.
    Supplied course document: \code{#2}.
  \end{thebibliography}}

\providecommand{\tightlist}{\setlength{\itemsep}{0pt}\setlength{\parskip}{0pt}}
\begin{document}
\makeutmreporttitlepage
\tableofcontents
\clearpage
\begin{utmabstract}[Abstract]
Two xv6 user programs demonstrate the process lifecycle and communication between a parent and child. Execution experiments compare successful program replacement with a missing executable. A one-byte exchange over two pipes demonstrates blocking reads, descriptor ownership and deterministic message order.
\utmkeywords{xv6, fork, exec, wait, pipe, interprocess communication}
\end{utmabstract}
\utmintroduction

Two xv6 user programs exercise process creation and pipe communication.
\texttt{run} starts a child that becomes \texttt{echo}; additional runs
replace it with \texttt{ls} and attempt to execute a nonexistent
program. \texttt{pingpong} exchanges one byte between a parent and child
using two pipes.

The programs were built against xv6 revision
\texttt{06aad25c735fd3159bdfae5680be4eab7b1668b2} using the RISC-V GCC
toolchain and executed in QEMU with one virtual CPU. Both the kernel and
filesystem image built successfully. Complete build and execution output
is retained in \texttt{evidence/}.

The submitted \texttt{Makefile} adds the following tab-indented entries
to \texttt{UPROGS}:

\begin{lstlisting}
    $U/_run\
    $U/_pingpong\
\end{lstlisting}

The corresponding sources are \texttt{user/run.c} and
\texttt{user/pingpong.c}. The existing \texttt{args} and \texttt{sleep}
entries are retained, with their sources included so that the submitted
Makefile has all its additional user programs. Copy the Makefile and
submitted user sources into the matching xv6 tree, then build and boot
with \texttt{make\ qemu}. Exit QEMU with Ctrl-A followed by X before
rebuilding a changed program.

\chapter{Process creation inside xv6}\label{process-creation-inside-xv6}

\section{Running echo}\label{running-echo}

\texttt{run} prints its PID and calls \texttt{fork()}. The child prints
its own PID and constructs a null-terminated argument vector:

\begin{lstlisting}
char *args[] = { "echo", "hello", "from", "the", "child", 0 };
exec("echo", args);
\end{lstlisting}

The parent calls \texttt{wait(\&status)} and prints the collected exit
status. The observed output was:

\begin{lstlisting}
$ run
parent: my pid is 3
child: my pid is 4
hello from the child
parent: child 4 ended with status 0
\end{lstlisting}

The parent and child have different PIDs. After successful
\texttt{exec}, the child remains process 4 but executes \texttt{echo}
rather than the instructions of \texttt{run}. The final parent message
appears after \texttt{echo} terminates.

\section{Replacing echo with ls}\label{replacing-echo-with-ls}

The child argument vector and execution call were changed to:

\begin{lstlisting}
char *args[] = { "ls", 0 };
exec("ls", args);
\end{lstlisting}

After rebuilding and rebooting, selected output was:

\begin{lstlisting}
$ run
parent: my pid is 3
child: my pid is 4
.              1 1 1024
..             1 1 1024
README         2 2 2441
run            2 3 34584
pingpong       2 4 35720
[remaining directory entries omitted]
parent: child 4 ended with status 0
\end{lstlisting}

The complete listing is in \texttt{evidence/ls.txt}. The directory
listing confirms that the child executed \texttt{ls}, and its normal
termination produced status 0. PIDs repeat because these variants were
tested in separate boots.

\newpage
\section{Failed execution}\label{failed-execution}

Changing the execution call to \texttt{exec("nosuchprog",\ args);},
rebuilding and running again produced:

\begin{lstlisting}
$ run
parent: my pid is 3
child: my pid is 4
exec failed
parent: child 4 ended with status 1
\end{lstlisting}

The error message comes from the code immediately after \texttt{exec}.
That code is reachable only when execution fails. The child then calls
\texttt{exit(1)}. The submitted source retains the successful
\texttt{echo} version.

\section{Comparison with Linux:
observations}\label{comparison-with-linux-observations}

\begin{enumerate}
\def\labelenumi{\arabic{enumi}.}
\tightlist
\item
  xv6 uses \texttt{fork()}, \texttt{exec("echo",\ args)} and
  \texttt{wait(\&status)}. The Linux program uses \texttt{fork()},
  \texttt{execlp("ls",\ "ls",\ "-l",\ (char\ *)NULL)} and
  \texttt{waitpid(pid,\ \&status,\ 0)}. Both follow the same sequence:
  create a child, replace its program, and collect its termination
  status. xv6's \texttt{exec} takes an argument array and does not
  perform the Linux \texttt{execlp} search through \texttt{PATH}; its
  \texttt{wait} reaps an available child rather than selecting a PID.
  There is only one child here, so that distinction does not affect the
  experiment.
\item
  The code of \texttt{echo} prints ``hello from the child''.
  \texttt{run} supplies those words as arguments; after successful
  \texttt{exec}, its child-side code has been replaced.
\item
  The failed-exec status is 1 in xv6 and 127 in the Linux experiment
  because the programs explicitly call \texttt{exit(1)} and
  \texttt{exit(127)}, respectively. These are program choices, not
  different automatic translations by the kernels. xv6 returns the
  integer passed to \texttt{exit} through \texttt{wait}; Linux uses an
  encoded wait status from which the program extracts
  \texttt{WEXITSTATUS} for a normal exit.
\end{enumerate}

\chapter{Bidirectional communication with
pipes}\label{bidirectional-communication-with-pipes}

\section{Protocol and descriptor
ownership}\label{protocol-and-descriptor-ownership}

\texttt{pingpong} creates \texttt{p2c}, for parent-to-child traffic, and
\texttt{c2p}, for child-to-parent traffic, before calling \texttt{fork}.
The parent initializes the byte to \texttt{x}. The child receives it,
prints its PID and ``received ping'', sends the same byte back, closes
its pipe ends and exits. The parent receives and checks the returned
byte, prints its PID and ``received pong'', and waits for the child.

{
\begin{longtable}[]{@{}llll@{}}
\hline
Process & Reads from & Writes to & Unused ends closed immediately \\
\hline
\endhead
\hline
\endlastfoot
Parent & \texttt{c2p{[}0{]}} & \texttt{p2c{[}1{]}} &
\texttt{p2c{[}0{]}}, \texttt{c2p{[}1{]}} \\
Child & \texttt{p2c{[}0{]}} & \texttt{c2p{[}1{]}} & \texttt{p2c{[}1{]}},
\texttt{c2p{[}0{]}} \\
\end{longtable}
}

The essential operations are \texttt{read(fd,\ \&b,\ 1)} and
\texttt{write(fd,\ \&b,\ 1)}. Their return values must equal one for a
successful one-byte exchange. The implementation checks pipe creation,
\texttt{fork}, reads, writes and the child's exit status. It closes each
used end after its final operation. On an error exit, xv6 also closes
the process's remaining descriptors.

The ordering is enforced by data dependencies. The child prints before
sending the reply, and the parent prints only after reading that reply.
No sleep is required.

\section{Observed exchanges}\label{observed-exchanges}

Three consecutive executions in the same boot produced:

\begin{lstlisting}
$ pingpong
6: received ping
5: received pong
$ pingpong
8: received ping
7: received pong
$ pingpong
10: received ping
9: received pong
\end{lstlisting}

Every execution printed exactly one ping followed by one pong, then
returned to the shell. The ping PID belongs to the child; the pong PID
belongs to its parent. The increasing values reflect newly created
processes, not a fixed requirement of the protocol.

\newpage
\section{Pipe observations}\label{pipe-observations}

\begin{enumerate}
\def\labelenumi{\arabic{enumi}.}
\tightlist
\item
  Two pipes provide two separate one-way channels. A pipe is a shared
  byte stream, not an addressed message connection. If both processes
  read and write the same pipe without an additional synchronization
  protocol, the parent can consume the byte it just sent instead of a
  reply. Separate pipes make the sender and receiver unambiguous in each
  direction.
\item
  \texttt{read(c2p{[}0{]},\ \&b,\ 1)} makes the parent wait for the
  reply when the return pipe is empty and a writer remains open. The
  later \texttt{wait(\&status)} reaps the child; it is not the operation
  that waits for the reply byte. If the child has already written the
  byte before the parent calls \texttt{read}, that read can complete
  immediately.
\end{enumerate}

\utmconclusions

The xv6 experiments reproduce the same process lifecycle used on Linux
with a smaller user API. Successful \texttt{exec} preserves process
identity while changing the executing program, and failed execution
returns to code that chooses an exit status. The ping-pong program
demonstrates that blocking reads synchronize processes through data
availability. Separate pipes, deliberate descriptor ownership and a
final wait provide an ordered exchange and proper child cleanup.

\end{document}
